\chapter{Calendarizacion y secuenciacion}

\section{Definicion}
La calendarizacion asigna trabajos a recursos en el tiempo. En manufactura y
servicios, sus objetivos frecuentes son minimizar el makespan, tardanza,
tiempos ociosos o costos por incumplimiento \cite{pinedo2016scheduling}.

\section{Flow shop y job shop}
En un flow shop, todos los trabajos visitan las maquinas en el mismo orden. En
un job shop, cada trabajo puede tener una ruta propia. Esta diferencia altera
profundamente la dificultad del problema y el tipo de representacion requerida.

\section{Estandar para ejemplos}
Los ejemplos del manual deberan incluir tabla de operaciones, diagrama de Gantt,
makespan, secuencia de trabajos y salida JSON de \qsb{} con tiempos de inicio y
fin por operacion.
